
\documentclass[12pt,a4paper]{article}

\usepackage[margin=1in]{geometry}
\usepackage{setspace}
\usepackage{parskip}
\usepackage{enumitem}
\usepackage{amsmath,amssymb}
\usepackage{siunitx}
\usepackage{microtype}
\usepackage{hyperref}
\hypersetup{hidelinks}
\setstretch{1.12}

\begin{document}

\begin{center}
{\LARGE\bfseries Novelty and Contribution Statement}

\vspace{0.7em}
{\large\itshape A Frame-Invariant Evaluation Protocol for Cross-Platform Short-Horizon UAV Inertial Forecasting}

\vspace{0.4em}
Pradip Diwan Borase and Vijyant Agarwal\\
Department of Mechanical Engineering, Netaji Subhas University of Technology, New Delhi, India
\end{center}

\section*{1. Novelty in one sentence}

The novelty of this work is \textbf{not} the Euclidean norm, a new neural architecture, or IMU forecasting by itself; it is a \textbf{physically defensible and leakage-controlled cross-platform evaluation protocol} that enables short-horizon forecasting of a common UAV inertial quantity across heterogeneous public datasets without fitting a target-domain axis alignment, while preserving source-only model selection, causal preprocessing, flight-level inference, and frozen zero-shot transfer.

\section*{2. Technical gap addressed}

Cross-dataset evaluation of learning-based inertial forecasting is difficult because nominally similar IMU channels may differ in sensor-frame orientation, native sampling rate, platform dynamics, recording duration, and data-processing conventions. A comparison based directly on native $x$-, $y$-, and $z$-components can therefore conflate forecasting performance with an unsupported coordinate-frame assumption. In addition, target-domain normalization, alignment, model selection, or early stopping can turn a nominally ``zero-shot'' comparison into implicit adaptation, while treating thousands of overlapping forecast windows as independent samples can substantially overstate the effective experimental sample size.

The present study addresses these issues using EuRoC MAV and UZH-FPV. It defines accelerometer specific-force magnitude
\[
q_f(t)=\|\mathbf{f}(t)\|_2
\]
as the primary forecast quantity and angular-speed magnitude
\[
q_\omega(t)=\|\boldsymbol{\omega}(t)\|_2
\]
as an auxiliary input. The Euclidean norm is standard mathematics and is explicitly \textbf{not} claimed as novel. Its role is to provide a physically transparent scalar quantity whose value is unchanged by a fixed orthonormal rotation of the sensor axes, thereby avoiding an unsupported target-fitted EuRoC--UZH body-axis registration.

\section*{3. Principal methodological contributions}

The paper makes three linked methodological contributions.

\begin{enumerate}[leftmargin=*]
    \item \textbf{Frame-invariant cross-platform quantity definition.}
    A common scalar inertial quantity is defined for cross-dataset comparison when a justified common body-axis transformation is unavailable. This preserves the physical unit of specific force while making the evaluation invariant to a fixed sensor-axis rotation. The limitation is stated explicitly: directional information is discarded, so the formulation is not a substitute for vector-valued state estimation or equivariant inertial odometry.

    \item \textbf{Leakage-controlled evaluation protocol.}
    The protocol combines causal anti-alias filtering and deterministic resampling to \SI{100}{\hertz}, a common \SI{1.0}{\second} candidate input window, a \SI{0.20}{\second} direct forecast horizon, source-only scaling and model selection, frozen source-to-target transfer, and flight recording---not overlapping forecast window---as the inferential unit. All target-domain fitting, alignment, early stopping, and architecture selection are prohibited.

    \item \textbf{Complexity--generalization benchmark under physically heterogeneous transfer.}
    Persistence, source-selected Ridge, TCN, GRU, and a compact Transformer are evaluated under the same forecasting task. The central empirical contribution is that added temporal-model complexity gives \emph{domain-dependent} gains rather than uniformly superior cross-platform generalization. This conclusion is strengthened by paired flight-level analysis, complete 10--200~ms horizon profiles, and sequence-level bootstrap intervals.
\end{enumerate}

\section*{4. Fairness and reproducibility refinements}

Two technical audits materially strengthen the reproducibility of the final manuscript.

\paragraph{Filter-order audit.}
The actual preprocessing code uses SciPy elliptic IIR design in SOS form with a \SI{30}{\hertz} passband edge, \SI{45}{\hertz} stopband edge, \SI{1}{\decibel} maximum passband ripple, and \SI{60}{\decibel} minimum stopband attenuation. The resulting \SI{200}{\hertz} EuRoC filter is \textbf{fifth order}, while the \SI{500}{\hertz} UZH-FPV filter is \textbf{sixth order}. The earlier ``order 6'' EuRoC label arose from inferring order as twice the number of SOS rows; the executed coefficients and processed signals were not changed. The manuscript now reports the actual executed orders and preserves the exact coefficients.

\paragraph{Input-history fairness audit.}
Every prediction origin provides a common 100-sample (\SI{1.0}{\second}) candidate history. TCN, GRU, and Transformer consume all 100 samples. Ridge selects an effective suffix length from the predeclared grid $\{10,25,50,100\}$ using source-only validation. The primary B3 Ridge therefore uses 25 samples for EuRoC and 10 samples for UZH-FPV. To test whether this difference in effective memory unfairly favors Ridge, a post hoc equal-history sensitivity fixes Ridge at 100 samples and retunes only $\alpha$ under the original source-only rules. The fixed-100-lag RMSE is worse than the primary source-selected-lag Ridge in all four contexts:
\begin{itemize}[leftmargin=*]
    \item Within EuRoC: 0.748 versus 0.728~\si{\metre\per\second\squared};
    \item Within UZH-FPV: 1.957 versus 1.943~\si{\metre\per\second\squared};
    \item EuRoC$\rightarrow$UZH-FPV: 2.126 versus 2.065~\si{\metre\per\second\squared};
    \item UZH-FPV$\rightarrow$EuRoC: 0.842 versus 0.831~\si{\metre\per\second\squared}.
\end{itemize}
No neural model is retrained, and target-domain data are not used to tune this sensitivity analysis. The result shows that Ridge was not denied access to the same one-second history; source validation simply preferred shorter effective memory.

\section*{5. Main empirical finding}

The study does not support a universal ``deep learning is better'' conclusion. Within EuRoC, source-selected Ridge is strongest on the primary aggregate metric. Within UZH-FPV, the Transformer provides the clearest nonlinear gain. Under EuRoC$\rightarrow$UZH transfer, that Transformer advantage reverses, while TCN remains close to Ridge. Under UZH-FPV$\rightarrow$EuRoC transfer, GRU performs best. Across the four contexts, the best learned method improves RMSE relative to Persistence by approximately 12.2--20.6\%, confirming nontrivial short-horizon predictability despite the gravity-dominated scalar target.

The scientific message is therefore that \textbf{within-dataset model ranking is not a reliable proxy for cross-platform generalization}, and that evaluation design---physical quantity definition, causal preprocessing, leakage control, inferential unit, and frozen transfer---can be as important as architecture choice.

\section*{6. What is explicitly not claimed}

The manuscript does not claim:
\begin{itemize}[leftmargin=*]
    \item a novel Euclidean-norm operation;
    \item a new TCN, GRU, Transformer, or Ridge architecture;
    \item that nonlinear models universally outperform Ridge;
    \item that domain shift is eliminated;
    \item that $q_f$ is independently referenced translational acceleration;
    \item that the reported bootstrap intervals are a GUM-style metrological uncertainty budget;
    \item onboard real-time flight suitability; or
    \item generalization to all UAV platforms.
\end{itemize}

\section*{7. Significance for \emph{Measurement}}

The work is positioned as a measurement-oriented evaluation study rather than as an architecture-proposal paper. Its contribution is a reproducible procedure for comparing learning-assisted inertial forecasting across heterogeneous sensing platforms while respecting coordinate-frame ambiguity, causal signal processing, source-only fitting, and the true independent experimental unit. The study therefore contributes to reliable benchmarking and interpretation of machine-learning-assisted inertial measurements, particularly where apparent performance improvements may otherwise arise from frame assumptions, target leakage, unequal recording influence, or model-selection bias.

\end{document}
